\section{Kết quả đạt được}

\subsection{Tối ưu mô hình YOLO}

Đồ án đã xây dựng pipeline tối ưu YOLO26n sáu giai đoạn, mỗi giai đoạn là một
script độc lập và được đóng gói thành notebook chạy lại được trên Google Colab:
\begin{itemize}
  \item Hiện thực thuật toán Sea-Horse Optimizer độc lập với framework và dùng nó
        để tìm đồng thời 18 siêu tham số tối ưu hóa và tăng cường dữ liệu.
  \item Mở rộng Ultralytics với các module cắt tỉa cho kiến trúc YOLO26
        (\texttt{C3k2Pruned}, \texttt{C2PSAPruned}, \texttt{SPPFPruned},
        \texttt{DetectPruned}, \texttt{DetectionModelPruned}), huấn luyện thưa hóa
        BatchNorm với hệ số giảm dần và cắt tỉa theo ngưỡng toàn cục có giới hạn an
        toàn. Cắt tỉa 30\% giảm số tham số từ 2,57 triệu xuống 2,17 triệu ($-15{,}7\%$)
        và FLOPs từ 6,12 xuống 4,82 GFLOPs ($-21{,}2\%$).
  \item Hiện thực chưng cất tri thức Channel-Wise Distillation kết hợp chưng cất
        điểm phân lớp cho mô hình đã cắt tỉa.
  \item Hiện thực QAT cho YOLO26 đã cắt tỉa bằng \texttt{pytorch-quantization},
        gồm các lớp lượng tử tùy chỉnh cho phép cộng dư/ghép kênh và hiệu chỉnh
        entropy, cho phép xuất ONNX Q/DQ sang TensorRT INT8.
  \item Các thực nghiệm sơ bộ rút ra ba nhận xét thực tiễn: lịch giảm dần của hệ
        số thưa hóa là cần thiết; mô hình sau cắt cần ngân sách phục hồi đủ dài;
        trọng số chưng cất lớn gây hại khi huấn luyện ngắn và độ biến thiên theo
        seed đủ lớn để đòi hỏi so sánh nhiều lần chạy.
\end{itemize}

\subsection{Nền tảng Edge AI}

Đồ án đã xây dựng và triển khai một nền tảng Cloud--Edge hoàn chỉnh:
\begin{itemize}
  \item \textbf{Cloud Control Plane} một tiến trình phục vụ REST, WebSocket, MCP
        và bộ tiêu thụ MQTT; quản lý cluster, node, runtime, deployment; lưu
        PostgreSQL và Redis; theo dõi liveness của node.
  \item \textbf{Edge Agent} chạy đa kiến trúc (x86, ARM64), điều khiển container
        qua Docker theo mô hình desired state: thao tác idempotent, chống lệnh
        trùng, lưu desired state xuống đĩa, tiếp quản container sau khi khởi động
        lại, vòng reconcile tự tạo lại container lỗi với giới hạn số lần thử.
  \item \textbf{AI Runtime} chạy YOLO (PyTorch, ONNX hoặc TensorRT) trên luồng
        camera với hàng đợi giới hạn, tự đo hơn 25 chỉ số runtime, camera và đầu ra
        model, tự suy ra trạng thái DEGRADED/FROZEN.
  \item \textbf{Phân phối model} qua kho Hugging Face Hub ghim theo commit, tải về
        Edge có cache và ghi nguyên tử.
  \item \textbf{Giám sát} hai kênh: telemetry ứng dụng theo đặc tả thống nhất và
        metric hạ tầng Prometheus (agent mode, remote write) với luật cảnh báo;
        thiết kế mô-đun cảnh báo sớm dựa trên đặc trưng cửa sổ trượt.
  \item \textbf{Web Dashboard} React theo dõi realtime qua WebSocket.
  \item \textbf{Đóng gói và CI/CD}: các stack Docker Compose, năm workflow GitHub
        Actions build image lên GHCR và triển khai tự động lên AWS EC2.
\end{itemize}
Bộ 18 kiểm thử tự động chạy mã sản phẩm của Cloud và Agent qua bus trong bộ nhớ
đạt 17/18, bao phủ luồng triển khai đầu cuối, tự phục hồi, telemetry, realtime,
xác thực và công cụ MCP.

\subsection{AI Assistant}

AI Assistant được xây dựng trên LangGraph với đồ thị bốn nút
(react--response--genchart--suggestion), truy cập hệ thống duy nhất qua 17 công
cụ MCP do Cloud công bố. Công cụ chỉ định nghĩa một lần ở Cloud và dùng chung
service với REST API; công cụ ghi bị ẩn theo mặc định và chỉ được gọi sau khi
người dùng xác nhận; lỗi công cụ được trả về mô hình như dữ liệu. Câu trả lời
được phát trực tuyến qua SSE kèm trạng thái gọi công cụ, biểu đồ và gợi ý câu hỏi.
